\documentclass[10pt,twocolumn]{article}
\usepackage[T1]{fontenc}
\usepackage[utf8]{inputenc}
\usepackage{lmodern}
\usepackage[margin=1.8cm,top=2.4cm,bottom=2cm,columnsep=0.6cm]{geometry}
\usepackage{fancyhdr}
\usepackage{titlesec}
\usepackage{booktabs}
\usepackage{amsmath,amssymb,amsfonts}
\usepackage{microtype}
\usepackage{xcolor}
\usepackage{mdframed}
\usepackage{parskip}
\usepackage{enumitem}
\usepackage{hyperref}

\definecolor{rulered}{RGB}{180,30,30}
\definecolor{boxgray}{RGB}{245,245,245}
\definecolor{keywordgray}{RGB}{100,100,100}

\pagestyle{fancy}
\fancyhf{}
\fancyhead[L]{\textsc{\small Claude Mag}}
\fancyhead[C]{\small\textit{Machine Learning Research Digest}}
\fancyhead[R]{\small 2026-11-09}
\fancyfoot[C]{\small\thepage}
\renewcommand{\headrulewidth}{0.8pt}

\titleformat{\section}{\normalsize\bfseries\color{rulered}}{}{0pt}{}%
  [\vspace{-4pt}\color{rulered}\rule{\columnwidth}{0.4pt}]
\titlespacing{\section}{0pt}{8pt}{4pt}

\hypersetup{colorlinks=true,urlcolor=rulered,linkcolor=black}

\begin{document}
\setlength{\parindent}{0pt}
\setlength{\parskip}{4pt}

{\small\color{keywordgray}\textbf{Keywords:}
  LLM inference \textbullet{}
  quantization \textbullet{}
  memory hierarchy \textbullet{}
  KV cache}\par\vspace{4pt}

{\LARGE\bfseries\raggedright
  StrataFlux: Memory-Hierarchy-Native Low-Bit LLM Inference}\par\vspace{4pt}

{\large\itshape\raggedright
  Using Host DRAM and NVMe SSD as Precision Strata Recovers Near-Full
  Accuracy at 2-Bit Weight Quantization Without Extra GPU Memory}\par\vspace{6pt}

{\small\textbf{By} Jake Hyun, Ziyu Zhao, Jae-sun Seo, Mohamed S.\ Abdelfattah}\par\vspace{2pt}

{\small\color{keywordgray}arXiv:2611.08363 \textbullet{}
  cs.LG / cs.DC \textbullet{} submitted 5 Nov 2026}
\par\vspace{2pt}\noindent{\color{rulered}\rule{\columnwidth}{0.6pt}}\vspace{6pt}

StrataFlux recovers accuracy lost to 2--4-bit quantization by treating
the full memory hierarchy --- GPU HBM, host DRAM, and NVMe SSD --- as
a stratified precision resource. On 2-bit Llama-3.1-8B-Instruct,
WikiText-2 perplexity drops from 9.46 to 8.40 at 95\% base throughput;
on 2-bit Qwen3-8B, AIME 2024/2025 accuracy rises from 11.0\% to 47.1\%
with no additional GPU memory.

\begin{mdframed}[backgroundcolor=boxgray,linecolor=rulered,linewidth=1pt,
    innerleftmargin=6pt,innerrightmargin=6pt,
    innertopmargin=6pt,innerbottommargin=6pt]
\textbf{\small AT A GLANCE}\par\vspace{4pt}
\begin{description}[leftmargin=0pt,labelsep=4pt,itemsep=2pt,parsep=0pt,
    font=\small\bfseries]
  \item[Problem:] 2-bit quantization drastically reduces GPU memory
    but collapses accuracy; existing mixed-precision fixes keep all
    higher-precision weights on-GPU, negating the memory benefit.
  \item[Why it matters:] Enabling 2-bit inference at acceptable accuracy
    would halve peak GPU memory for many deployments, making 70B+
    models accessible on single-GPU servers.
  \item[Method:] Stratified sensitivity profiling assigns weight layers
    to GPU HBM (2-bit), DRAM (4-bit), or NVMe SSD (8-bit); a prefetch
    scheduler hides tier-B/C transfer latency behind GPU compute.
  \item[Result:] 2-bit Llama-3.1-8B perplexity 8.40 (vs.\ 9.46
    baseline); 2-bit Qwen3-8B AIME accuracy 47.1\% (vs.\ 11.0\%).
  \item[Evidence:] WikiText-2, AIME 2024/2025, 128K long-context;
    compared against 2-bit RTN, GPTQ, AWQ; RTX 5090.
\end{description}
\end{mdframed}

\section{The Big Picture}

Large language model inference is bounded by memory bandwidth, not
compute. Weights must be loaded from GPU memory on every forward pass,
and peak GPU memory sets the maximum model size that fits on a single
device. Quantization reduces both: 2-bit weights need 16$\times$ less
memory than FP32 and consume less bandwidth per layer.

The problem is accuracy. At 2 bits, the quantization error per weight
is large enough to damage the model's reasoning ability, particularly
on tasks that require precise recall of numerical facts or long
multi-step chains. Prior solutions (GPTQ, AWQ, SpQR) identify sensitive
weights and store them at higher precision \emph{on the GPU}, but this
cancels the memory benefit for the sensitive fraction.

StrataFlux reframes the question: the GPU is not the only memory in a
modern server. Host DRAM is 8--32$\times$ larger and 5--10$\times$
slower; NVMe SSD is 10--100$\times$ larger still. The sensitive weights
that need 4- or 8-bit precision can live in DRAM or SSD and be streamed
to GPU only when needed for a specific layer, overlapped with the GPU
computation on the previous layer.

\section{Why Existing Approaches Fall Short}

\textbf{Uniform quantization} (RTN, 2-bit) is fast and memory-efficient
but loses too much accuracy in sensitive layers, especially attention
projections and early embedding layers.

\textbf{Mixed-precision on-GPU} (GPTQ, AWQ, SpQR) keeps sensitive
weights at 4--8 bits on the GPU. This recovers accuracy but requires
extra GPU memory proportional to the number of sensitive parameters,
partially defeating the purpose of quantization.

\textbf{CPU offloading} (FlexGen, DejaVu) moves entire layers to CPU
for latency-insensitive serving but introduces large latency spikes when
layers are paged back in, making throughput unpredictable.

StrataFlux differs from CPU offloading in two ways: (1) it offloads
only the \emph{higher-precision copies} of sensitive layers, not the
full layer; the base 2-bit copy always stays on GPU for fast path; and
(2) a prefetch scheduler computes the access pattern ahead of time and
hides DMA transfers behind GPU computation.

\section{Core Mechanism: Three-Tier Precision}

\textbf{Stratum A (GPU HBM):} The base quantized weights of all layers
at the target bit-width (2--4 bits). Always resident; used for fast-path
inference.

\textbf{Stratum B (host DRAM):} 4-bit copies of the top-$k$\% most
sensitive layers, identified by calibration. Paged in one block ahead
of the layer it corrects.

\textbf{Stratum C (NVMe SSD):} 8-bit copies of an extremely sensitive
subset (typically the first and last 2 transformer blocks, plus any
layer where calibration error exceeds a threshold). Paged in
asynchronously via direct DMA.

For a layer assigned to stratum B or C, the forward pass combines:
\[
  \mathbf{y} = W_{\text{base}} \mathbf{x}
    + \Delta W_{\text{tier}} \mathbf{x},
\]
where $W_{\text{base}}$ is the 2-bit on-GPU weight and
$\Delta W_{\text{tier}} = W_{\text{tier}} - W_{\text{base}}$ is the
correction loaded from DRAM or SSD. This residual form ensures that
if the DMA transfer has not completed (e.g., SSD latency spike), the
base path produces an approximate result instead of blocking.

\section{Technical Formulation}

Sensitivity of layer $l$ is estimated by the Fisher trace:
\[
  S_l = \mathrm{tr}\bigl(\mathbf{H}_l\bigr)
    \approx \frac{1}{|C|} \sum_{x \in C}
    \bigl\|W_l x - \hat{W}_l x\bigr\|_2^2,
\]
where $C$ is a calibration corpus, $W_l$ is the full-precision weight,
and $\hat{W}_l$ is its 2-bit quantized version. Layers are ranked by
$S_l$ and assigned to strata by two thresholds $\theta_B$ and
$\theta_C$:
\begin{align*}
  l \in \text{Stratum A} &\;\text{if}\; S_l \le \theta_B,\\
  l \in \text{Stratum B} &\;\text{if}\; \theta_B < S_l \le \theta_C,\\
  l \in \text{Stratum C} &\;\text{if}\; S_l > \theta_C.
\end{align*}

Thresholds $\theta_B$ and $\theta_C$ are set by a memory budget
constraint: the DRAM and SSD footprints are fixed by the available
hardware, and layers are sorted by $S_l$ until the budget is consumed.

\section{Learning or Inference Procedure}

StrataFlux requires no fine-tuning; it modifies only the inference
memory layout and access pattern. The workflow is:
\begin{enumerate}[itemsep=2pt,parsep=0pt]
  \item Run calibration pass on 128 calibration sequences to compute
    $S_l$ for all layers.
  \item Assign layers to strata based on thresholds.
  \item Quantize and place stratum-B copies in DRAM, stratum-C in SSD.
  \item At inference time, the prefetch scheduler issues DMA requests
    for the next 3 stratum-B layers and 1 stratum-C layer ahead of
    execution, overlapping transfers with GPU computation.
\end{enumerate}

\section{What the Guarantee Says}

No theoretical accuracy bound is provided. The empirical claim is
Pareto-dominance over existing mixed-precision methods on the two
primary benchmarks: StrataFlux achieves lower perplexity at lower GPU
memory than any configuration of GPTQ or AWQ tested. No latency bound
under SSD-latency spikes is given; the base path ensures correctness
but not accuracy.

\section{Experimental Findings}

\begin{itemize}[itemsep=2pt,parsep=0pt]
  \item \textbf{Perplexity:} 2-bit Llama-3.1-8B-Instruct WikiText-2
    PPL $8.40$ (StrataFlux) vs.\ $9.46$ (2-bit RTN base); 95\% base
    throughput on RTX 5090.
  \item \textbf{Reasoning:} 2-bit Qwen3-8B AIME 2024/2025 accuracy
    $47.1\%$ (StrataFlux) vs.\ $11.0\%$ (2-bit RTN); GPU memory
    unchanged.
  \item \textbf{Long-context:} 3-bit weights + 2-bit KV cache at 128K
    context matches vLLM 4-bit AWQ + 8-bit KV perplexity and speed
    at half the GPU memory.
  \item \textbf{Throughput:} Prefetch scheduler achieves 95\% of base
    2-bit throughput by hiding DRAM latency; SSD accesses visible only
    when DRAM budget is exhausted.
\end{itemize}

\section{Ablations and Interpretation}

\textbf{Stratum-C necessity:} Removing stratum C (SSD tier) and
assigning all stratum-C layers to stratum B (DRAM at 4-bit) increases
DRAM footprint by 18\% at no accuracy change, confirming that the
SSD tier covers only a small but critical fraction of sensitive layers.

\textbf{Residual form:} Replacing the residual correction
$\Delta W_{\text{tier}}$ with a full replacement (always load
$W_{\text{tier}}$ on this path) reduces throughput by 12\% because
the base 2-bit GPU path is always bypassed, even when DMA is on time.

\textbf{Calibration size:} 32 calibration sequences achieve 94\% of the
sensitivity ordering quality of 512 sequences; 128 sequences reach 99\%.

\textbf{Model size:} Relative perplexity improvement is larger for
8B-class models (1.1 PPL points) than for 70B-class models (0.4 PPL
points), reflecting that smaller models are more sensitive to 2-bit
quantization and benefit more from stratum correction.

\section*{Reference}

Jake Hyun, Ziyu Zhao, Jae-sun Seo, Mohamed S.\ Abdelfattah.
\textbf{StrataFlux: Memory-Hierarchy-Native Low-Bit LLM Inference.}
arXiv:2611.08363.
\url{https://arxiv.org/abs/2611.08363}

\end{document}
